% !TeX root = ../main.tex

\chapter{Conclusion and Future Work}\label{chapter:conclusion}

\section{Conclusion}

This thesis demonstrates that video encoding can supply the geometric
observations needed for useful visual odometry. Motion estimation performed for
compression can also support camera tracking without a separate pixel-based
correspondence frontend. The successful trajectories, most clearly on the easy
Machine Hall sequences, establish that this reuse is feasible. Their significance
is that a representation produced for video coding can serve as the measurement
source for geometric estimation.

This opens a route to repurposing widely available video-encoding accelerators
for the odometry component of a SLAM system. The demonstrated method is naturally
compatible with compress-then-analyze at its analysis interface: a capturing
device can provide the prepared coded streams, and a receiver can estimate motion
from their codec metadata and calibration. The receiving device need not repeat
correspondence estimation on image intensities. For a small camera paired with a
phone, this offers a way to build localisation around existing coding resources
and the communication structure already used for compressed video.

The pipeline also permits other choices of where to place the communication
boundary. Its filtering and geometric validation stages expose intermediate
observations that could be transmitted instead of the video streams. A capturing
device could extract motion vectors, perform some of that analysis locally,
and then compress and send the retained correspondences for further estimation.
Such a division resembles analyze-then-compress at the transmission boundary,
while still obtaining its initial measurements through video compression. It
therefore need not be assigned exclusively to either CTA or ATC: the same
codec-based analysis principle supports a range of allocations between camera
and receiver. The useful split would depend on local computation, link capacity,
and the information required by the remaining estimation stages.

The contribution is a feasibility demonstration for this design direction.
The present software-encoded experiments establish useful geometric information;
their accuracy limits leave room for different correspondence constructions and
estimators. Hardware execution and intermediate-data transport remain systems
work to evaluate. What the thesis establishes is a basis for that work: visual
odometry can be built from codec-generated motion information, making existing
video-processing resources relevant to localisation as well as communication.

\section{Future Work}

A first priority is a causal hardware implementation that connects a constrained
capturing device to a receiving phone. The current two-pass temporal encoding
needs a streaming alternative. Offline reverse-temporal streams must be
replaced by bounded buffering or another causal mechanism, and
the shifted stereo inputs must be scheduled within the encoder's throughput
and latency limits. A mosaic construction, fewer pre-shifts, or another means
of obtaining depth could reduce input-preparation cost. Access to motion-vector
metadata should also be measured directly, including any image reconstruction
performed inside the extraction path.

The resulting system should be evaluated for energy per processed stereo pair,
power, latency, transmission rate, and trajectory accuracy at the same operating
point. Measurements should include rectification, every additional encode,
memory movement, the radio link, metadata extraction, and estimation on both
devices. Comparing image-based CTA, codec-based CTA, and capture-side feature
extraction would test when the proposed allocation is beneficial. This follows
the processing--communication trade-off demonstrated by Redondi et
al.~\cite{redondi2016cta}, rather than transferring a hardware power figure to
a different complete workload. Sending an ordinary coded stream and sending a
dedicated motion-vector payload are separate transport options to evaluate.

A second priority is to improve the geometric observations while preserving
the pixel-free estimation boundary. The stereo examples show how grid phases
refine spatial sampling and how pre-shifts can make depth available at different
locations. Future work could select these transforms adaptively, compare hardware
encoder motion fields, and account more explicitly for block support and
measurement uncertainty. Observation models and association rules should be
judged by their ability to retain geometrically valid connections as well as
by their trajectory error. A longer assigned track alone does not establish
that it follows the same scene point.

The preliminary experiments discussed in
\autoref{appendix:prior-experiments} also motivate maintaining surface or
region identities while their supporting codec correspondences change.
Such representations could reduce dependence on long-lived point tracks,
provided that their association and geometry constrain camera motion.

Registered RGB-D measurements offer a complementary experiment. Replacing the
codec-stereo stage with sensor depth would help separate temporal motion-vector
limitations from missing stereo support, while keeping temporal tracking free
of image-intensity access. This would be a different information setting from
the stereo system studied here, and would test whether the same analysis
principle is useful with another depth source.

Finally, evaluation should extend to other cameras, encoders, and motion
patterns, including controlled exploration and mapping routes. Such tasks can
make moderate, repeatable motion an operational choice, motivating study of
where the favourable tracking seen here is useful. Tests under rapid rotation,
blur, low texture, and changing illumination would establish the limits of that
use. The present method provides a starting point for these investigations:
codec motion can sustain useful camera tracking, and efficient systems built
around that information are worth exploring further.
